\documentclass[UTF8,11pt]{ctexart}
\usepackage[a4paper,margin=22mm]{geometry}
\usepackage{amsmath,amssymb,bm,graphicx,adjustbox,fvextra,longtable,array}
\xeCJKDeclareCharClass{CJK}{"2160 -> "217F, "2460 -> "249B, "2200 -> "22FF}
\newcommand{\fitmath}[1]{\begin{adjustbox}{max width=\linewidth}$\displaystyle #1$\end{adjustbox}}
\setlength{\parindent}{0pt}\setlength{\parskip}{.65em}\renewcommand{\arraystretch}{1.3}
\sloppy\allowdisplaybreaks
\begin{document}
\section*{2004 年考研政治真题}
\subsection*{2004 年硕士研究生入学考试政治试题}

\subsection*{一、单项选择题}

1. 唯物史观认为，人类的第一个历史活动是


\begin{center}\begin{tabular}{>{\raggedright\arraybackslash}p{\dimexpr(\linewidth-8\tabcolsep)/4\relax}>{\raggedright\arraybackslash}p{\dimexpr(\linewidth-8\tabcolsep)/4\relax}>{\raggedright\arraybackslash}p{\dimexpr(\linewidth-8\tabcolsep)/4\relax}>{\raggedright\arraybackslash}p{\dimexpr(\linewidth-8\tabcolsep)/4\relax}}
\textbf{A.} 吃喝穿住 & \textbf{B.} 物质生活资料的生产 & \textbf{C.} 人的自觉意识活动 & \textbf{D.} 结成社会关系\\[.35em]\end{tabular}\end{center}

2. 20 世纪50 年代，北大荒人烟稀少、一片荒凉。由于人口剧增，生产力水平低下，吃饭问题成为中国面临的首要问题，于是人们不得不靠扩大耕地面积增加粮食产量，经过半个世纪的开垦，北大荒成了全国闻名的“北大仓”。然而由于过度开垦已经造成了许多生态问题。现在，黑龙江垦区全面停止开荒，退耕还“荒”。这说明


\begin{center}\begin{tabular}{>{\raggedright\arraybackslash}p{\dimexpr(\linewidth-4\tabcolsep)/2\relax}>{\raggedright\arraybackslash}p{\dimexpr(\linewidth-4\tabcolsep)/2\relax}}
\textbf{A.} 人与自然的和谐最终以恢复原始生态为归宿 & \textbf{B.} 人们改造自然的一切行为都会遭到“自然界的报复”\\[.35em]
\textbf{C.} 人在自然界面前总是处于被支配的地位 & \textbf{D.} 人们应合理地调节人与自然之间的物质变换\\[.35em]\end{tabular}\end{center}

3.在抗击“非典”的斗争中，许多患者被治愈后又捐出自己的血清，用于治疗其他患者,这说明


\begin{center}\begin{tabular}{>{\raggedright\arraybackslash}p{\dimexpr(\linewidth-4\tabcolsep)/2\relax}>{\raggedright\arraybackslash}p{\dimexpr(\linewidth-4\tabcolsep)/2\relax}}
\textbf{A.} 人的价值只体现在特定的场合和行为中 & \textbf{B.} 人的价值必须以满足个人需要为前提\\[.35em]
\textbf{C.} 人的价值是在满足自身和他人的需要中实现的 & \textbf{D.} 人的价值表现了人的能力的大小\\[.35em]\end{tabular}\end{center}

4. 社会总资本扩大再生产的前提条件是


\begin{center}\begin{tabular}{>{\raggedright\arraybackslash}p{\dimexpr(\linewidth-8\tabcolsep)/4\relax}>{\raggedright\arraybackslash}p{\dimexpr(\linewidth-8\tabcolsep)/4\relax}>{\raggedright\arraybackslash}p{\dimexpr(\linewidth-8\tabcolsep)/4\relax}>{\raggedright\arraybackslash}p{\dimexpr(\linewidth-8\tabcolsep)/4\relax}}
\textbf{A.} Ⅰ(V＋M)＝ⅡC & \textbf{B.} Ⅱ(V＋M)＝ⅠC & \textbf{C.} Ⅰ(V＋M)>ⅡC & \textbf{D.} Ⅱ(V＋M)>ⅠC\\[.35em]\end{tabular}\end{center}

5. 把公司全部资本分为等额股份，股东以其出资额为限对公司承担责任，公司以其全部资产对公司的债务承担责任，这是


\begin{center}\begin{tabular}{>{\raggedright\arraybackslash}p{\dimexpr(\linewidth-8\tabcolsep)/4\relax}>{\raggedright\arraybackslash}p{\dimexpr(\linewidth-8\tabcolsep)/4\relax}>{\raggedright\arraybackslash}p{\dimexpr(\linewidth-8\tabcolsep)/4\relax}>{\raggedright\arraybackslash}p{\dimexpr(\linewidth-8\tabcolsep)/4\relax}}
\textbf{A.} 无限责任公司 & \textbf{B.} 股份有限公司 & \textbf{C.} 有限责任公司 & \textbf{D.} 合伙制企业\\[.35em]\end{tabular}\end{center}

6. 新民主主义革命的中心内容是


\begin{center}\begin{tabular}{>{\raggedright\arraybackslash}p{\dimexpr(\linewidth-4\tabcolsep)/2\relax}>{\raggedright\arraybackslash}p{\dimexpr(\linewidth-4\tabcolsep)/2\relax}}
\textbf{A.} 没收封建地主阶级的土地归新民主主义国家所有 & \textbf{B.} 没收官僚垄断资本归新民主主义国家所有\\[.35em]
\textbf{C.} 没收封建地主阶级的土地归农民所有 & \textbf{D.} 保护民族工商业\\[.35em]\end{tabular}\end{center}

7. 毛泽东首次明确提出“新民主主义革命”这一科学概念的著作是


\begin{center}\begin{tabular}{>{\raggedright\arraybackslash}p{\dimexpr(\linewidth-4\tabcolsep)/2\relax}>{\raggedright\arraybackslash}p{\dimexpr(\linewidth-4\tabcolsep)/2\relax}}
\textbf{A.} 《〈共产党人〉发刊词》 & \textbf{B.} 《中国革命和中国共产党》\\[.35em]
\textbf{C.} 《新民主主义论》 & \textbf{D.} 《论联合政府》\\[.35em]\end{tabular}\end{center}

8.在中国共产党七届二中全会上，毛泽东告诫全党：“务必使同志们继续地保持谦虚、谨慎、不骄、不躁的作风，务必使同志们继续地保持艰苦奋斗的作风。”其原因主要是


\begin{center}\begin{tabular}{>{\raggedright\arraybackslash}p{\dimexpr(\linewidth-4\tabcolsep)/2\relax}>{\raggedright\arraybackslash}p{\dimexpr(\linewidth-4\tabcolsep)/2\relax}}
\textbf{A.} 中国共产党即将成为执政党 & \textbf{B.} 党的工作方式发生了变化\\[.35em]
\textbf{C.} 全国大陆即将解放 & \textbf{D.} 中国将由新民主主义社会转变为社会主义社会\\[.35em]\end{tabular}\end{center}

9. 在农业社会主义改造中建立的初级农业生产合作社属于


\begin{center}\begin{tabular}{>{\raggedright\arraybackslash}p{\dimexpr(\linewidth-8\tabcolsep)/4\relax}>{\raggedright\arraybackslash}p{\dimexpr(\linewidth-8\tabcolsep)/4\relax}>{\raggedright\arraybackslash}p{\dimexpr(\linewidth-8\tabcolsep)/4\relax}>{\raggedright\arraybackslash}p{\dimexpr(\linewidth-8\tabcolsep)/4\relax}}
\textbf{A.} 新民主主义性质 & \textbf{B.} 社会主义萌芽性质 & \textbf{C.} 半社会主义性质 & \textbf{D.} 社会主义性\\[.35em]\end{tabular}\end{center}

10. 贯彻“三个代表”重要思想，关键在


\begin{center}\begin{tabular}{>{\raggedright\arraybackslash}p{\dimexpr(\linewidth-8\tabcolsep)/4\relax}>{\raggedright\arraybackslash}p{\dimexpr(\linewidth-8\tabcolsep)/4\relax}>{\raggedright\arraybackslash}p{\dimexpr(\linewidth-8\tabcolsep)/4\relax}>{\raggedright\arraybackslash}p{\dimexpr(\linewidth-8\tabcolsep)/4\relax}}
\textbf{A.} 坚持党的先进性 & \textbf{B.} 坚持执政为民 & \textbf{C.} 坚持党的阶级性 & \textbf{D.} 坚持与时俱进\\[.35em]\end{tabular}\end{center}

11. “三个代表”重要思想是我们党的立党之本、执政之基、力量之源。这里的“本”、“基”、“源”，说到底就是


\begin{center}\begin{tabular}{>{\raggedright\arraybackslash}p{\dimexpr(\linewidth-4\tabcolsep)/2\relax}>{\raggedright\arraybackslash}p{\dimexpr(\linewidth-4\tabcolsep)/2\relax}}
\textbf{A.} 发展先进生产力 & \textbf{B.} 发展先进文化\\[.35em]
\textbf{C.} 人民群众的支持和拥护 & \textbf{D.} 人民群众生活水平的提高\\[.35em]\end{tabular}\end{center}

12. 某员工在外资企业工作，年薪5 万元；利用业余时间在民营企业兼职，年薪2 万元；购买股票分得的红利2 万元；出租住房收入2 万元；转让一项技术收入1 万元。该员工一年的劳动收入为


\begin{center}\begin{tabular}{>{\raggedright\arraybackslash}p{\dimexpr(\linewidth-8\tabcolsep)/4\relax}>{\raggedright\arraybackslash}p{\dimexpr(\linewidth-8\tabcolsep)/4\relax}>{\raggedright\arraybackslash}p{\dimexpr(\linewidth-8\tabcolsep)/4\relax}>{\raggedright\arraybackslash}p{\dimexpr(\linewidth-8\tabcolsep)/4\relax}}
\textbf{A.} 12 万元 & \textbf{B.} 9 万元 & \textbf{C.} 8 万元 & \textbf{D.} 7 万元\\[.35em]\end{tabular}\end{center}

13. 2003 年9 月通过全民公决否决了加入欧元区议案的欧盟国家是


\begin{center}\begin{tabular}{>{\raggedright\arraybackslash}p{\dimexpr(\linewidth-8\tabcolsep)/4\relax}>{\raggedright\arraybackslash}p{\dimexpr(\linewidth-8\tabcolsep)/4\relax}>{\raggedright\arraybackslash}p{\dimexpr(\linewidth-8\tabcolsep)/4\relax}>{\raggedright\arraybackslash}p{\dimexpr(\linewidth-8\tabcolsep)/4\relax}}
\textbf{A.} 瑞典 & \textbf{B.} 丹麦 & \textbf{C.} 英国 & \textbf{D.} 希腊\\[.35em]\end{tabular}\end{center}

14. 2003 年月22 日至27 日，印度总理瓦杰帕伊对我国进行正式访问，这是印度总理10 年来首次访华，访问取得了成功，两国签署了


\begin{center}\begin{tabular}{>{\raggedright\arraybackslash}p{\dimexpr(\linewidth-4\tabcolsep)/2\relax}>{\raggedright\arraybackslash}p{\dimexpr(\linewidth-4\tabcolsep)/2\relax}}
\textbf{A.} 《中印关系原则和全面合作宣言》 & \textbf{B.} 《中印联合新闻公报》\\[.35em]
\textbf{C.} 《中印联合声明》 & \textbf{D.} 《关于中国西藏地方和印度之间的通商和交通协定》\\[.35em]\end{tabular}\end{center}

15. 2003 年中国与欧盟关系有了进一步发展，其突出表现是


\begin{center}\begin{tabular}{>{\raggedright\arraybackslash}p{\dimexpr(\linewidth-4\tabcolsep)/2\relax}>{\raggedright\arraybackslash}p{\dimexpr(\linewidth-4\tabcolsep)/2\relax}}
\textbf{A.} 中国与欧盟领导人年度会晤机制起步 & \textbf{B.} 中国与欧盟建立全面伙伴关系\\[.35em]
\textbf{C.} 中国首次制定和发表《中国对欧盟政策文件》 & \textbf{D.} 欧盟首次制定和发表《中国欧盟关系长期政策》\\[.35em]\end{tabular}\end{center}

\subsection*{二、多项选择题}

16. 2003 年6 月23 日，《城市生活无着的流浪乞讨人员救助管理办法》正式发布，并于8 月1 日正式实施。1982 年发布的《城市流浪乞讨人员收容遣送办法》同时被废止。这一变化体现了


\begin{center}\begin{tabular}{>{\raggedright\arraybackslash}p{\dimexpr(\linewidth-4\tabcolsep)/2\relax}>{\raggedright\arraybackslash}p{\dimexpr(\linewidth-4\tabcolsep)/2\relax}}
\textbf{A.} 政治文明的进步 & \textbf{B.} 对人民群众根本利益的维护\\[.35em]
\textbf{C.} 对人权的尊重和保护 & \textbf{D.} 上层建筑不断变革完善的要求\\[.35em]
\textbf{E.} 生产关系的根本变革 & \\[.35em]\end{tabular}\end{center}

17. 有一幅广告幽默画，画的是几个行人在看一家饭店外贴的告示，上写：“快进来吃饭吧，否则你我都得挨饿。”这幅广告画的寓意有


\begin{center}\begin{tabular}{>{\raggedright\arraybackslash}p{\dimexpr(\linewidth-4\tabcolsep)/2\relax}>{\raggedright\arraybackslash}p{\dimexpr(\linewidth-4\tabcolsep)/2\relax}}
\textbf{A.} 生产者和消费者是相互依存的 & \textbf{B.} 生产和消费具有直接的同一性\\[.35em]
\textbf{C.} 利己是人的一切活动的出发点 & \textbf{D.} 商品交换活动背后隐藏着人与人的关系\\[.35em]
\textbf{E.} 生产关系本质上是人与人之间的物质利益关系 & \\[.35em]\end{tabular}\end{center}

18. 社会主义市场经济中按劳分配的“劳”是指


\begin{center}\begin{tabular}{>{\raggedright\arraybackslash}p{\dimexpr(\linewidth-4\tabcolsep)/2\relax}>{\raggedright\arraybackslash}p{\dimexpr(\linewidth-4\tabcolsep)/2\relax}}
\textbf{A.} 不同形式的具体劳动 & \textbf{B.} 通过价值形式来实现的劳动\\[.35em]
\textbf{C.} 直接的社会劳动 & \textbf{D.} 公有制企业的局部劳动\\[.35em]
\textbf{E.} 劳动者实际支出的劳动 & \\[.35em]\end{tabular}\end{center}

19. 利润转化为平均利润的过程，同时也是


\begin{center}\begin{tabular}{>{\raggedright\arraybackslash}p{\dimexpr(\linewidth-4\tabcolsep)/2\relax}>{\raggedright\arraybackslash}p{\dimexpr(\linewidth-4\tabcolsep)/2\relax}}
\textbf{A.} 资本有机构成提高的过程 & \textbf{B.} 价值转化为生产价格的过程\\[.35em]
\textbf{C.} 资本在不同部门之间发生转移的过程 & \textbf{D.} 资本家集团重新瓜分剩余价值的过程\\[.35em]
\textbf{E.} 超额利润消失的过程 & \\[.35em]\end{tabular}\end{center}

20. 新民主主义革命和社会主义革命的关系是


\begin{center}\begin{tabular}{>{\raggedright\arraybackslash}p{\dimexpr(\linewidth-4\tabcolsep)/2\relax}>{\raggedright\arraybackslash}p{\dimexpr(\linewidth-4\tabcolsep)/2\relax}}
\textbf{A.} 新民主主义革命与社会主义革命可以“毕其功于一役” & \textbf{B.} 新民主主义革命和社会主义革命的任务相同\\[.35em]
\textbf{C.} 两个革命之间需要有一个资本主义的过渡阶段 & \textbf{D.} 新民主主义革命是社会主义革命的必要准备\\[.35em]
\textbf{E.} 社会主义革命是新民主主义革命的必然趋势 & \\[.35em]\end{tabular}\end{center}

21. 中国共产党在把马克思列宁主义基本原理与中国革命实际相结合的过程中，在学风问题上曾经反对过的主要错误倾向是


\begin{center}\begin{tabular}{>{\raggedright\arraybackslash}p{\dimexpr(\linewidth-4\tabcolsep)/2\relax}>{\raggedright\arraybackslash}p{\dimexpr(\linewidth-4\tabcolsep)/2\relax}}
\textbf{A.} 投降主义 & \textbf{B.} 经验主义\\[.35em]
\textbf{C.} 教条主义 & \textbf{D.} 冒险主义\\[.35em]
\textbf{E.} 机会主义 & \\[.35em]\end{tabular}\end{center}

22. 在社会主义改造基本完成以后，正确处理人民内部矛盾成为国家政治生活的主题。中国共产党提出的正确处理人民内部矛盾的方针政策主要有


\begin{center}\begin{tabular}{>{\raggedright\arraybackslash}p{\dimexpr(\linewidth-4\tabcolsep)/2\relax}>{\raggedright\arraybackslash}p{\dimexpr(\linewidth-4\tabcolsep)/2\relax}}
\textbf{A.} 统筹兼顾、适当安排 & \textbf{B.} 有理、有利、有节\\[.35em]
\textbf{C.} 百花齐放、百家争鸣 & \textbf{D.} 长期共存、相互监督\\[.35em]
\textbf{E.} 和平共处五项原则 & \\[.35em]\end{tabular}\end{center}

23. 1959 年底至1960 年初，毛泽东在读苏联《政治经济学教科书》时，认为社会主义社会的发展阶段有


\begin{center}\begin{tabular}{>{\raggedright\arraybackslash}p{\dimexpr(\linewidth-4\tabcolsep)/2\relax}>{\raggedright\arraybackslash}p{\dimexpr(\linewidth-4\tabcolsep)/2\relax}}
\textbf{A.} 不发达的社会主义 & \textbf{B.} 比较发达的社会主义\\[.35em]
\textbf{C.} 发达的社会主义 & \textbf{D.} 社会主义初级阶段\\[.35em]
\textbf{E.} 社会主义高级阶段 & \\[.35em]\end{tabular}\end{center}

24. 社会主义精神文明在社会主义现代化建设中具有的重要战略地位和作用，表现在


\begin{center}\begin{tabular}{>{\raggedright\arraybackslash}p{\dimexpr(\linewidth-4\tabcolsep)/2\relax}>{\raggedright\arraybackslash}p{\dimexpr(\linewidth-4\tabcolsep)/2\relax}}
\textbf{A.} 为经济建设提供精神动力 & \textbf{B.} 为现代化建设提供智力支持\\[.35em]
\textbf{C.} 为现代化建设创造良好的社会环境 & \textbf{D.} 为现代化建设的正确发展方向提供思想保证\\[.35em]
\textbf{E.} 它是构成综合国力的重要方面 & \\[.35em]\end{tabular}\end{center}

25. 邓小平关于社会主义本质的论断体现了


\begin{center}\begin{tabular}{>{\raggedright\arraybackslash}p{\dimexpr(\linewidth-4\tabcolsep)/2\relax}>{\raggedright\arraybackslash}p{\dimexpr(\linewidth-4\tabcolsep)/2\relax}}
\textbf{A.} 解放生产力与发展生产力的统一 & \textbf{B.} 生产力与生产关系的统一\\[.35em]
\textbf{C.} 发展生产力与实现共同富裕的统一 & \textbf{D.} 目的与手段的统一\\[.35em]
\textbf{E.} 社会主义发展过程与最终目标的统一 & \\[.35em]\end{tabular}\end{center}

26. 我国新型工业化道路的内涵是


\begin{center}\begin{tabular}{>{\raggedright\arraybackslash}p{\dimexpr(\linewidth-4\tabcolsep)/2\relax}>{\raggedright\arraybackslash}p{\dimexpr(\linewidth-4\tabcolsep)/2\relax}}
\textbf{A.} 科技含量高 & \textbf{B.} 经济效益好\\[.35em]
\textbf{C.} 资源消耗低 & \textbf{D.} 产值增加快\\[.35em]
\textbf{E.} 环境污染少 & \\[.35em]\end{tabular}\end{center}

27. “股份制是现代企业的一种资本组织形式，不能笼统地说股份制是公有还是私有”，这一观点表明


\begin{center}\begin{tabular}{>{\raggedright\arraybackslash}p{\dimexpr(\linewidth-4\tabcolsep)/2\relax}>{\raggedright\arraybackslash}p{\dimexpr(\linewidth-4\tabcolsep)/2\relax}}
\textbf{A.} 由法人股东而不是个人股东构成的股份制是公有制 & \textbf{B.} 公有制与私有制都可以通过股份制这一形式来实现\\[.35em]
\textbf{C.} 有公有制经济参股的就是公有制 & \textbf{D.} 股份制本身不具有公有还是私有的性质\\[.35em]
\textbf{E.} 公有制经济占控股地位就具有明显的公有性 & \\[.35em]\end{tabular}\end{center}

28. “中东和平路线图”计划实施进程困难重重，主要原因有


\begin{center}\begin{tabular}{>{\raggedright\arraybackslash}p{\dimexpr(\linewidth-4\tabcolsep)/2\relax}>{\raggedright\arraybackslash}p{\dimexpr(\linewidth-4\tabcolsep)/2\relax}}
\textbf{A.} 以色列右翼势力的阻拦 & \textbf{B.} 巴勒斯坦激进组织的干扰\\[.35em]
\textbf{C.} 美国偏袒以色列 & \textbf{D.} 伊拉克战争的爆发\\[.35em]
\textbf{E.} 俄罗斯对“中东和平路线图”计划持反对态度 & \\[.35em]\end{tabular}\end{center}

29. 2003 年8 月27-29 日在北京举行的“六方会谈”，为和平解决朝鲜核问题迈出了重要一步，其主要收获有


\begin{center}\begin{tabular}{>{\raggedright\arraybackslash}p{\dimexpr(\linewidth-4\tabcolsep)/2\relax}>{\raggedright\arraybackslash}p{\dimexpr(\linewidth-4\tabcolsep)/2\relax}}
\textbf{A.} 确定了朝鲜半岛无核化的目标，同时也都认识到需要考虑和解决朝鲜在安全等方面提出的关切 & \textbf{B.} 确认了通过对话以和平方式解决朝鲜半岛核问题的途径，维护半岛的和平与稳定\\[.35em]
\textbf{C.} 原则赞同解决核问题采取分阶段、同步或并行实施的方针，探讨并确定公正合理的总体解决方案 & \textbf{D.} 主张保持对话、建立信任、减少分歧、扩大共识，在和谈进程中不采取可能使局势升级或激化的言行\\[.35em]
\textbf{E.} 美朝双方初步达成了签署互不侵犯条约的共识 & \\[.35em]\end{tabular}\end{center}

30. 中共十六届三中全会强调，完善社会主义市场经济体制的主要任务是


\begin{center}\begin{tabular}{>{\raggedright\arraybackslash}p{\dimexpr(\linewidth-4\tabcolsep)/2\relax}>{\raggedright\arraybackslash}p{\dimexpr(\linewidth-4\tabcolsep)/2\relax}}
\textbf{A.} 完善公有制为主体、多种所有制经济共同发展的基本经济制度 & \textbf{B.} 建立有利于逐步改变城乡二元经济结构的体制，形成促进区域经济协调发展的机制\\[.35em]
\textbf{C.} 建设统一开放竞争有序的现代市场体系，完善宏观调控体系、行政管理体制和经济法律制度 & \textbf{D.} 健全就业、收入分配和社会保障制度\\[.35em]
\textbf{E.} 建立促进经济社会可持续发展的机制 & \\[.35em]\end{tabular}\end{center}

\subsection*{三、辨析题}

31. 用马克思主义哲学有关原理对漫画中所反映的工作方式进行辨析。


32. 2002 年9 月初，国家食品药品监管局等七部门联手出击开展屠宰市场集中整治，严把肉品市场准入关，规定进入市场销售的肉品必须是由定点屠宰厂(场)生产、经检疫合格的产品。对失信企业实行“黑名单”制度，其违法违规行为将被记录在案，公开曝光。根据材料辨析：有必要限制市场机制的作用，让拥有强制力的政府来干预。


33. 在新民主主义革命时期，只有当民族资产阶级拥护革命时，才要保护民族资本主义。


\subsection*{四、分析题}

34. 闻一多有一次给学生上课，他走上讲台，先在黑板上写了一道算术题：2＋5＝？学生们疑惑不解。然而闻先生却执意要问：2＋5＝？同学们于是回答：“等于7 嘛！”闻先生说：“不错。在数学领域里2＋5＝7，这是天经地义的颠扑不破的。但是，在艺术领域里，2＋5＝10000 也是可能的。”他拿出一幅题为《万里驰骋》的图画叫学生们欣赏，只见画面上突出地画了两匹奔马，在这两匹奔马后面，又错落有致、大小不一地画了五匹马，这五匹马后面便是许多影影绰绰的黑点点了。闻先生指着画说：“从整个画面的形象看，只有前后七匹马，然而，凡是看过这幅画的人，都会感到这里有万马奔腾，这难道不是2＋5＝10000 吗？”运用认识论相关原理分析下列问题：


（1）既然在数学领域2+5=7 是颠扑不破的，为什么在艺术领域2+5=10000 也是可能的？


（2）在认识活动中，正确处理理性与非理性的关系对科学创新有何重要意义？


35. 据估计，今天在美国有6000 家公司推行“雇员拥有股票计划”，其中包括西尔斯－罗伯克百货公司、美国电话电报公司等。“雇员拥有股票计划”在这些公司的推行，使工人们积极地经营他们的公司，产生了一种充满活力的责任感，在生产率、高质量和低成本等方面取得了巨大的成就。美国争取雇员拥有股票全国委员会对350 家高技术公司所作的一项调查发现，利用雇员拥有股票计划的公司要比没有利用这种计划的公司发展快2～4 倍。随着这一计划的推行，到2000 年，全美国有25％的雇员分享他们公司的所有权。这种迅速出现的“工人资本主义”概念也适用于相当大部分的美国经济。但是工人拥有股票不会轻易转变为工人管理。有的工人股东说：“我看不出有什么变化。一切都和以前一模一样。”也有的工人股东认为，在“雇员拥有股票计划”下，越是尽力干，得到的就越多。摘自W. E. 哈拉尔著：《新资本主义》请回答：


（1）根据材料分析当代资本主义社会实行“雇员拥有股票计划”的原因。


（2）评析工人股东的两种看法。


\begin{center}\begin{adjustbox}{max width=\linewidth}\begin{tabular}{|>{\raggedright\arraybackslash}p{\dimexpr(\linewidth-12\tabcolsep-7\arrayrulewidth)/6\relax}|>{\raggedright\arraybackslash}p{\dimexpr(\linewidth-12\tabcolsep-7\arrayrulewidth)/6\relax}|>{\raggedright\arraybackslash}p{\dimexpr(\linewidth-12\tabcolsep-7\arrayrulewidth)/6\relax}|>{\raggedright\arraybackslash}p{\dimexpr(\linewidth-12\tabcolsep-7\arrayrulewidth)/6\relax}|>{\raggedright\arraybackslash}p{\dimexpr(\linewidth-12\tabcolsep-7\arrayrulewidth)/6\relax}|>{\raggedright\arraybackslash}p{\dimexpr(\linewidth-12\tabcolsep-7\arrayrulewidth)/6\relax}|}\hline
 & 农业在国民 & 农业劳动力占 & 城镇化 & 城镇居民人均 & 农民人均\\\hline
 & 经济中的比重 & 总就业的比重 & 比例 & 可支配收入 & 可支配收入\\\hline
2000 年 & 15.\allowbreak{}9\% & 50\% & 36.\allowbreak{}2\% & 6280 元 & 2253 元\\\hline
2020 年 & 11.\allowbreak{}5\% & 29\% & 56\% & 18000 元 & 8000 元\\\hline\end{tabular}\end{adjustbox}\end{center}

（1）请结合材料说明为什么解决好“三农”问题是全面建设小康社会的重大任务。


（2）对比表中农业劳动力占总就业比重的数字，你认为应采取哪些措施来降低农业劳动力的比重？


37. 本题为选做题，请在Ⅰ、Ⅱ两道试题中选取其中一道作答，若两题都回答，只按第Ⅰ道试题的成绩记入总分。


\subsection*{选做题Ⅰ：}

\subsection*{材料1}

两个小时前，盟国空军部队开始向伊拉克和科威特的军事目标发动袭击。在我讲话的同时，这些袭击正在继续进行……国决议和经美国国会同意采取这次军事行动之前，联合国、美国和许多其他国家进行了历时几个月的不间断的和几乎是无休止的外交活动……派有部队的28 家已尽了一切应尽的努力来达成和平解决，没有别的选择，只能用武力把萨达姆赶出科威特……成功以后，我们将得到建立世界新秩序的真正机会。在那个新秩序中，可以依赖的联合国可以发挥其维护和平的作用，将联合国创始人的诺言和远见卓识付诸实现。摘自乔治·布什1991 年1 月6 日的电视讲话


\subsection*{材料2}

各位公民，此时美国和联盟部队已经开始了军事行动，这次行动的目的是解除伊拉克的武装、解放这个国家的人民、保卫世界免遭严重危险……美国人民和我们的朋友以及盟友不会听任一个以大规模杀伤性武器威胁和平的非法政权摆布。我们现在就要用我们的陆军、空军、海军、海岸警卫队和海军陆战队对付这种威胁，这样将来我们就不必用大批的消防员、警察和医生在我们的大街上应付这种威胁……对我们的国家和整个世界的威胁将被消除。我们将度过这个危险时刻，并继续推进和平。我们将捍卫我们的自由。我们还将把自由带给其他人。摘自乔治·Ｗ·布什2003 年3 月19 日的电视讲话指出上述两次讲话内容的主要异同并分析其原因。


\subsection*{选做题Ⅱ：}

\subsection*{材料1}

20 世纪80 年代末90 年代初，东欧剧变、苏联解体，社会主义遭受重大挫折。对此西方思想界的保守派纷纷著书重新审视西方“胜利”的历史原因和人类历史发展道路。美国学者弗朗西斯·福山的《历史的终结及最后之人》便是这种背景的产物。福山在书中提出：一个值得注意的共识这几年已在世界出现，因为自由民主已克服世袭君主制、法西斯与共产主义这类相对的意识形态。自由民主可能形成“人类意识形态进步的终点”与“人类统治的最后形态”，也构成了“历史的终结”。自由民主的“理念”已不能再改良了。最值得注意的发展是，在拉丁美洲和东欧、苏联、中东与亚洲，强固的政府都在这二十年间动摇了。自由民主目前已及于全球的不同地区与文化，成为惟一一贯的政治憧憬对象。


\subsection*{材料2}

改革开放20 多年来，中国社会主义现代化建设以前所未有的速度向前发展，综合国力不断增强，人民生活总体达到小康水平。根据世界银行的排名，2001 年中国的经济总量仅次于美国、日本、德国、英国、法国，位居世界第六位。对外贸易进出口总额在世界上的排位上升到第5 位，2002 年中国吸收外资超过美国，成为吸收外资最多的国家。在下图中，作为衡量一国总体经济实力主要指标的GDP 也真实地反映了中国近五年来的变化。


（1）结合材料评析资本主义是“人类统治的最后形态”的观点。


（2）用唯物史观评析资本主义“自由民主的‘理念’已不能再改良了”的观点。


（3）在两制并存的背景下，社会主义国家如何才能更快地实现自身的发展？

\end{document}
